% Generated by analysis/build-paper.mjs from dist/survey-data.js. Do not edit by hand.
\subsection{00: A simple command}
\textbf{Ground truth.} A successful command passes through six layers in order: input $\rightarrow$ capability check $\rightarrow$ agent/reasoning (the external agent interprets the sentence and picks the matching function from the predefined function set it was given as context) $\rightarrow$ call validation $\rightarrow$ loading (the chosen function is loaded on the ESP32) $\rightarrow$ execution \& feedback. The language interpretation happens in the external agent, not on the rover's board, which only runs the predefined function. No code is written or compiled at run time. Nothing is learned from other users.\par
\noindent\texttt{normal\_multi} (core). \textbf{Key:} The request is checked against what the rover is physically able to do; An AI agent interprets the sentence and picks the matching predefined function; That function call is checked before it is sent to the rover; The chosen function is loaded onto the rover's board.
\begin{itemize}
\item The rover's own board reads and interprets the English sentence directly $\rightarrow$ \textbf{WM6}. The ESP32 only runs a predefined function; the sentence is interpreted by the external agent.
\item The rover looks up how other users phrased similar commands $\rightarrow$ \textbf{LS}. Nothing is learned from other users; each command is handled by the agent and pipeline in this session.
\end{itemize}
\noindent\texttt{normal\_open}. \textbf{Model answer:} The command goes to an external AI agent which interprets it and picks the matching function from the predefined set; the call is checked (validated), the function is loaded on the rover's board, and then executed by the motors, with feedback returned. \textbf{Key ideas (keyword-matched hints):} An agent / AI interprets the command and picks a predefined function; It is checked / validated before use; It is loaded onto the board; The motors then execute it.\par
\noindent\texttt{normal\_grid}.
\begin{itemize}
\item The natural-language sentence is understood on the rover's board (the ESP32). $\rightarrow$ \textbf{False} (wrong answer tagged \textbf{WM6}). The board runs the program; the interpretation is done by the external agent.
\item Each command is matched to a function that already exists on the rover; no new program is written. $\rightarrow$ \textbf{True}. The agent picks an existing function from the set fixed at configuration time and supplies its arguments; nothing is generated or flashed, and the call is validated before it is loaded.
\item The rover is told what it can do through a registered list of its hardware capabilities. $\rightarrow$ \textbf{True}. The rover's capability manifest (motors, LED, wheel encoders; no camera, microphone, distance sensor or gripper) is registered when it connects.
\end{itemize}
\subsection{WM2: The rover did not move}
\textbf{Ground truth.} The failure was at the Call Validation stage. The agent picked a function that is not in the rover's function set (stopMotorz()); the catalogue lookup failed, the argument and safety checks were skipped, and nothing was loaded or executed. The motors were fine. In the opaque interface this looks like a generic error, indistinguishable from an input, loading or motor-feedback failure; the transparent interface names the layer.\par
\noindent\texttt{wm2\_stage} (core). \textbf{Key:} The function call the agent chose was rejected by checks before loading (call validation).
\begin{itemize}
\item Your command never reached the system (input) $\rightarrow$ \textbf{WM2}. You placed the failure at a different layer than the one that failed.
\item The system decided the rover can't do this (capability check) $\rightarrow$ \textbf{WM2}. The rover can move forward; this was not a capability refusal.
\item The AI agent couldn't interpret the request (agent / reasoning) $\rightarrow$ \textbf{WM2}. The agent responded; the function it chose was the problem.
\item The chosen function couldn't be loaded onto the rover (loading) $\rightarrow$ \textbf{WM2}. The call never got as far as loading.
\item The rover ran the function but no motion was confirmed (execution \& feedback) $\rightarrow$ \textbf{WM2}. Nothing was loaded or run, so this stage was never reached.
\end{itemize}
\noindent\texttt{wm2\_cause} (core). \textbf{Key:} The agent picked a function that is not in the rover's function set, so the call was rejected and the rover never received it.
\begin{itemize}
\item The motors or wheels are broken $\rightarrow$ \textbf{WM2}. Treating ``no motion'' as a single permanent hardware fault. The motors were never asked to move.
\item The rover understood but decided not to move $\rightarrow$ \textbf{WM6}. The rover does not decide; a software step upstream failed.
\item My wording was wrong and a different phrasing would have worked $\rightarrow$ \textbf{WM2}. The command was valid; the fault was in the function the agent chose.
\end{itemize}
\noindent\texttt{wm2\_grid}.
\begin{itemize}
\item If the rover does not move, the motors must be faulty. $\rightarrow$ \textbf{False} (wrong answer tagged \textbf{WM2}). Input, function selection, loading and motor feedback faults can all look like ``nothing happened''.
\item In this attempt the selected function was loaded onto the rover's board. $\rightarrow$ \textbf{False}. Validation failed first, so loading was blocked (R3: validate before load).
\item Different kinds of fault can look identical from the outside: the rover just doesn't move. $\rightarrow$ \textbf{True}. That is why the transparent interface names the failing layer.
\item Replacing the motors would fix this failure. $\rightarrow$ \textbf{False} (wrong answer tagged \textbf{WM2}). The motors were not involved.
\item The same command can succeed later once the underlying fault is removed. $\rightarrow$ \textbf{True}. The failure was specific to this function selection, not a permanent inability.
\end{itemize}
\noindent\texttt{wm2\_open}. \textbf{Model answer:} The agent picked a function that is not in the function set and the call failed validation, so nothing was loaded; the motors were fine. Next: retry after the fault is fixed. \textbf{Key ideas (keyword-matched hints):} The agent picked a wrong or non-existent function and the call failed validation; It was not loaded, so the rover never ran it; The motors / hardware were not at fault; Retry after fixing the cause.\par
\subsection{WM3: Find the red ball}
\textbf{Ground truth.} The registered capability manifest lists no camera (or microphone, distance sensor or gripper). Finding a red ball needs visual input, so the request was refused at the Capability check stage, before any function was selected or loaded. No phrasing and no smarter AI can compensate for missing hardware. In the opaque interface, the refusal is generic.\par
\noindent\texttt{wm3\_can} (core). \textbf{Key:} No --- it has no camera (or other sensor able to see), so no phrasing can make this work.
\begin{itemize}
\item Yes --- better wording would make it work $\rightarrow$ \textbf{WM3}. The limit is hardware, not wording.
\item Yes --- it would drive around until it bumped into the ball $\rightarrow$ \textbf{WM3}. Without a way to see the ball, it cannot tell it from anything else; the system also does not invent partial actions.
\item Yes --- after it has learned from more examples $\rightarrow$ \textbf{LS}. The rover does not learn new senses from examples.
\end{itemize}
\noindent\texttt{wm3\_grid}.
\begin{itemize}
\item What the rover can do is limited by the hardware registered on its board, not only by how capable the AI is. $\rightarrow$ \textbf{True}. The capability manifest is grounded in installed hardware (R1/R2).
\item A more capable AI model would let the rover see without a camera. $\rightarrow$ \textbf{False} (wrong answer tagged \textbf{WM3}). A better model cannot supply a sensor the rover lacks.
\item The system picked a movement function for this request, tried it, and failed to find the ball. $\rightarrow$ \textbf{False} (wrong answer tagged \textbf{WM3}). The request was refused at the capability check, before any function was selected.
\item The refusal means the rover is broken. $\rightarrow$ \textbf{False} (wrong answer tagged \textbf{WM2}). It is a correct refusal, not a fault.
\item This rover has no camera, microphone, distance sensor or gripper. $\rightarrow$ \textbf{True}. Its manifest lists motors, an LED and wheel encoders only.
\end{itemize}
\noindent\texttt{wm3\_open}. \textbf{Model answer:} The rover would need a camera (vision hardware) registered in its capability list, or a different device that has one. Rewording the command would not help. \textbf{Key ideas (keyword-matched hints):} Needs a camera / vision sensor / different hardware; Wording or a smarter AI alone would not fix it.\par
\subsection{WM4: ``Do that again''}
\textbf{Ground truth.} The first command succeeded and became the agent's active context. The context was then expired immediately before the follow-up, so ``that'' could no longer be resolved: the failure was at the Agent / Reasoning stage, no function was selected and nothing was loaded. The record of the earlier action lives in the external agent's working context, not in the rover; the rover's motors and memory were not at fault and it did not ignore you. Restating the full command works.\par
\noindent\texttt{wm4\_why} (core). \textbf{Key:} The earlier action was no longer in the agent's active context, so ``that'' could not be resolved.
\begin{itemize}
\item The rover forgot or ignored me --- it is unreliable $\rightarrow$ \textbf{WM4}. Context limits are not forgetting or defiance.
\item The motors lost power or failed after the first move $\rightarrow$ \textbf{WM2}. The motors were never asked to move again.
\item The rover learned that repeating the move is unsafe $\rightarrow$ \textbf{LS}. The rover did not learn or judge anything.
\end{itemize}
\noindent\texttt{wm4\_where} (core). \textbf{Key:} In the AI agent's working context (its short-term conversation memory).
\begin{itemize}
\item In the rover's own memory (its board) $\rightarrow$ \textbf{WM6}. The board just runs predefined functions; conversational context is held by the agent.
\item In a permanent log that the rover learns from across sessions $\rightarrow$ \textbf{LS}. Context is temporary and session-bound.
\item Nowhere --- the system can never refer back to earlier commands $\rightarrow$ \textbf{WM4}. It can, while the earlier action is still in active context.
\end{itemize}
\noindent\texttt{wm4\_grid}.
\begin{itemize}
\item Restating the whole command (``Move forward for 2 seconds'') would work. $\rightarrow$ \textbf{True}. A full command does not depend on context.
\item The follow-up failed because the motors had a fault. $\rightarrow$ \textbf{False} (wrong answer tagged \textbf{WM2}). Context expiry is not a motor fault.
\item ``Do that again'' only works while an earlier action is still in the agent's active context. $\rightarrow$ \textbf{True}. That is the dependency the follow-up exposed.
\item After this failure the rover has permanently lost the ability to repeat actions. $\rightarrow$ \textbf{False} (wrong answer tagged \textbf{WM4}). A new full command works immediately.
\item The rover chose to ignore the request. $\rightarrow$ \textbf{False} (wrong answer tagged \textbf{WM6}). No choice was involved; the agent had nothing to refer to.
\end{itemize}
\noindent\texttt{wm4\_open}. \textbf{Model answer:} The first command succeeded and was held in the agent's context; that context expired, so ``do that again'' had nothing to refer to. The rover was not at fault, and restating the full command would work. \textbf{Key ideas (keyword-matched hints):} Memory / context expired or was limited; It is the agent (not the rover or motors) that holds it; Restating the full command would fix it.\par
\subsection{WM6: Who does the thinking?}
\textbf{Ground truth.} The simulated external agent was unavailable (a connection timeout at the Agent / Reasoning stage). The rover's board stayed connected but cannot interpret a new English request by itself, so no new function was selected. Interpretation happens in the agent, not the robot; the rover did not decide anything. The remedy is restoring the agent connection, not repairing the rover.\par
\noindent\texttt{wm6\_where} (core). \textbf{Key:} In an external AI agent that picks the matching function, which is then loaded on the board.
\begin{itemize}
\item Inside the rover's board (ESP32) --- it understands the sentence itself $\rightarrow$ \textbf{WM6}. The board runs functions; it does not interpret the sentence.
\item In the motors or wheel controller $\rightarrow$ \textbf{WM6}. Motors only act on a loaded function.
\item In the web page I typed into $\rightarrow$ \textbf{WM6}. The page relays the command; the interpretation is done by the agent.
\end{itemize}
\noindent\texttt{wm6\_why} (core). \textbf{Key:} The connection to the external agent timed out, so no function was selected, though the board stayed connected.
\begin{itemize}
\item The rover understood but chose not to act $\rightarrow$ \textbf{WM6}. The rover has no such choice.
\item The rover hardware has failed $\rightarrow$ \textbf{WM2}. The board remained connected and healthy.
\item The rover refused because turning left is unsafe $\rightarrow$ \textbf{WM6}. The safety check was never reached.
\end{itemize}
\noindent\texttt{wm6\_grid}.
\begin{itemize}
\item A connected board can interpret any new English request on its own. $\rightarrow$ \textbf{False} (wrong answer tagged \textbf{WM6}). Without the agent, no function is selected.
\item Restoring the agent connection, not repairing the rover, is what was needed. $\rightarrow$ \textbf{True}. The rover was not at fault.
\item No function was selected in this attempt. $\rightarrow$ \textbf{True}. The failure happened before function selection.
\item The rover ``decided'' not to turn. $\rightarrow$ \textbf{False} (wrong answer tagged \textbf{WM6}). No decision was made by the rover.
\end{itemize}
\noindent\texttt{wm6\_open}. \textbf{Model answer:} An external agent interprets the sentence and picks the function; the board only runs what is loaded. If the agent is unreachable, a connected board has nothing new to run. \textbf{Key ideas (keyword-matched hints):} An external agent / AI interprets the sentence; The board only runs the loaded function; The agent connection failed / timed out.\par
\subsection{WM7: It used to work}
\textbf{Ground truth.} The simulated firmware update lowered the maximum motion duration from 5 s to 3 s. The identical 4-second command now fails the safety check at Call Validation, so it is not loaded. The rover did not wear out and did not learn; a rule changed. The transparent interface shows a change notice (R8); the opaque interface hides it, so the change looks like an unexplained regression.\par
\noindent\texttt{wm7\_why} (core). \textbf{Key:} An update changed the allowed motion limit, so the safety check now blocks the 4-second command.
\begin{itemize}
\item The rover has worn out or got worse with use $\rightarrow$ \textbf{WM7}. Nothing physical degraded; a limit changed.
\item The rover learned that 4 seconds is risky and now refuses $\rightarrow$ \textbf{LS}. The rover did not learn anything.
\item It is a random glitch; retrying will eventually work $\rightarrow$ \textbf{WM7}. The identical command will keep failing until the rule or command changes.
\end{itemize}
\noindent\texttt{wm7\_grid}.
\begin{itemize}
\item A rule in the system can change even though I did nothing differently. $\rightarrow$ \textbf{True}. A firmware/model update changes behaviour without any user action.
\item Sending the identical 4-second command again will eventually succeed. $\rightarrow$ \textbf{False} (wrong answer tagged \textbf{WM7}). The safety check will block it each time.
\item A shorter command (e.g. 2 seconds) would still be allowed. $\rightarrow$ \textbf{True}. The new limit is 3 seconds.
\item The 4-second command was loaded and then stopped part-way. $\rightarrow$ \textbf{False}. It was blocked at validation, before loading.
\end{itemize}
\noindent\texttt{wm7\_open}. \textbf{Model answer:} A firmware update lowered the safety limit (5 s $\rightarrow$ 3 s). You could find out from a change notice, version information or a validation message. \textbf{Key ideas (keyword-matched hints):} A system/firmware/version change (not the rover ``wearing out''); The safety check / validation blocked it; A notice / changelog / version info would help.\par
\subsection{LEARN: After the fix}
\textbf{Ground truth.} The rover did not learn from the failure. The cause (the injected fault) was removed, so the newly selected function call passed validation. RoverLens keeps context only inside the current session; nothing is adapted, remembered across sessions or improved by repeated use, and using it does not change the rover's capabilities.\par
\noindent\texttt{ls\_why} (core). \textbf{Key:} The cause of the failure was removed, so the new attempt passed.
\begin{itemize}
\item The rover learned from its earlier failure and adapted $\rightarrow$ \textbf{LS}. No learning takes place; the fault itself was removed.
\item The rover just needed a few attempts to warm up $\rightarrow$ \textbf{WM2}. The failure had a specific cause; repeated attempts alone would not have changed the outcome.
\end{itemize}
\noindent\texttt{ls\_grid}.
\begin{itemize}
\item The rover now remembers this failure and will change its future behaviour. $\rightarrow$ \textbf{False} (wrong answer tagged \textbf{LS}). It does not adapt from failures.
\item A new session starts without the memory of this session. $\rightarrow$ \textbf{True}. Context and logs live only in the current session.
\item The rover's skills improve the more I use it. $\rightarrow$ \textbf{False} (wrong answer tagged \textbf{LS}). Its capabilities are fixed by hardware and firmware.
\item Repeated successful use can change what hardware capabilities the rover has. $\rightarrow$ \textbf{False} (wrong answer tagged \textbf{LS}). Use never adds sensors or actuators.
\end{itemize}
\noindent\texttt{ls\_open}. \textbf{Model answer:} No. Retries succeed because the cause was fixed, not because the rover learned. Context exists only within the current session and capabilities are fixed by hardware and firmware. \textbf{Key ideas (keyword-matched hints):} It does not learn / adapt; Context/memory only lasts within the session; The fix (fault removal) explains the retry.\par
\subsection{WM5: Vignette: early success}
\textbf{Ground truth.} One command type working repeatedly shows that the layers used by that task worked. Other tasks can use different layers or need capabilities that are missing (a camera), and agent availability, context and firmware can change. Blanket trust is miscalibrated, but so is blanket distrust: trust should be tied to task, capability and system state.\par
\noindent\texttt{wm5\_why} (core). \textbf{Key:} Not necessarily --- success at one command says little about other commands, which depend on different capabilities and system state.
\begin{itemize}
\item Yes --- five successes show the robot is reliable $\rightarrow$ \textbf{WM5}. Success on one task says little about a different task.
\item No --- robots like this are never reliable $\rightarrow$ \textbf{WM5}. Over-correction: the robot is reliable for tasks within its capabilities and healthy state.
\end{itemize}
\noindent\texttt{wm5\_grid}.
\begin{itemize}
\item After five successes, the chance of any failure is effectively zero. $\rightarrow$ \textbf{False} (wrong answer tagged \textbf{WM5}). Repeated success does not remove other failure sources.
\item A failure after several successes means the robot has become unreliable overall. $\rightarrow$ \textbf{False} (wrong answer tagged \textbf{WM5}). It may be a missing capability, an expired context, an agent outage or a change.
\item A task can fail because it needs a capability the robot never had, however many other tasks succeeded. $\rightarrow$ \textbf{True}. Camera-dependent tasks need a camera.
\end{itemize}
\subsection{WM1: Vignette: two languages}
\textbf{Ground truth.} The hardware did not change. The language interpretation and function selection are done by the external agent, which can perform unevenly across languages. A difference in results across languages points at the agent layer and its language coverage, not at broken hardware or the robot ``getting dumber''.\par
\noindent\texttt{wm1\_why} (core). \textbf{Key:} The agent that interprets language and picks the function may handle some languages less reliably, while the same hardware works fine.
\begin{itemize}
\item The robot's hardware is failing $\rightarrow$ \textbf{WM1}. The hardware is identical in both attempts.
\item The rover's board only understands English because English is built into it $\rightarrow$ \textbf{WM6}. The board doesn't interpret language at all.
\item The robot has gotten worse over time $\rightarrow$ \textbf{WM1}. Nothing degraded; the difference is in language coverage.
\end{itemize}
\noindent\texttt{wm1\_grid}.
\begin{itemize}
\item The failure necessarily means the motors changed. $\rightarrow$ \textbf{False} (wrong answer tagged \textbf{WM1}). Same hardware, different input language.
\item The same hardware can work well in one language and poorly in another. $\rightarrow$ \textbf{True}. Performance differences come from the language/agent layer.
\end{itemize}
